\documentclass[11pt,letterpaper]{article}
\usepackage[margin=1in]{geometry}
\usepackage[T1]{fontenc}
\usepackage{lmodern}
\usepackage{amsmath,amssymb,amsthm,mathtools}
\usepackage{microtype,booktabs,array,enumitem,graphicx}
\usepackage[dvipsnames]{xcolor}
\usepackage{fancyhdr}
\usepackage{xurl}
\usepackage{tocloft}
\setlength{\cftbeforesecskip}{2pt}
\usepackage[colorlinks=true,linkcolor=MidnightBlue,citecolor=MidnightBlue,urlcolor=MidnightBlue]{hyperref}
\newtheorem{theorem}{Theorem}[section]
\newtheorem{corollary}[theorem]{Corollary}
\newtheorem{proposition}[theorem]{Proposition}
\newtheorem{lemma}[theorem]{Lemma}
\theoremstyle{definition}
\newtheorem{definition}[theorem]{Definition}
\theoremstyle{remark}
\newtheorem{remark}[theorem]{Remark}
\newcommand{\Qp}{\mathbb Q_p}
\newcommand{\Zp}{\mathbb Z_p}
\newcommand{\Fp}{\mathbb F_p}
\setlength{\parskip}{4pt}
\setlength{\parindent}{0pt}
\setlength{\emergencystretch}{2em}
\pagestyle{fancy}
\fancyhf{}
\fancyhead[L]{\small ProveIt research extension}
\fancyhead[R]{\small November 27, 2026}
\fancyfoot[C]{\thepage}
\setlength{\headheight}{14pt}


\hypersetup{pdftitle={Quadratic rigidity of positive Fourier filters},pdfauthor={Research prepared for Vladimir Reshetnikov}}
\title{Quadratic Rigidity of Positive Fourier Filters\\\large The Hessian at every strictly positive mirror spectrum}
\author{Research prepared for Vladimir Reshetnikov}
\date{November 27, 2026}
\begin{document}
\maketitle
\begin{abstract}
We extend the source article's local positive-filter analysis from a
coalesced spectrum to every distinct reflection-paired real spectrum
whose optimal atomic weights are strictly positive. The width Hessian is
an explicit weighted variance integral. Its kernel is exactly the tangent
space preserving reflection symmetry; every transverse direction incurs
a positive quadratic width cost. A three-frequency specialization gives
a closed curvature formula throughout the positive band, including its
coalescence asymptotic and divergence at the vanishing-weight boundary.
This is an unconditional structural result answering part of an explicit
question in the source article. It is not a resolution of a major
classical conjecture, and publication priority remains unverified.
\end{abstract}
{\setlength{\parskip}{0pt}\tableofcontents}
\section{The transverse Hessian at an arbitrary positive mirror spectrum}

The repository article \cite{ProveItFilters} proves analytic continuation of the optimum at each strictly
positive reflection-paired configuration, and computes the width Hessian at a
coalesced frequency. Its second open question asks for the Hessian at a
general mirror configuration. The theorem below answers its definiteness and
kernel questions for distinct real frequencies (equivalently, exponents with
one common real part). It does not settle signs of all mixed minors, vanishing
base weights, repeated-frequency strata, or independently moving real parts.

For a list $\lambda=(\lambda_1,\ldots,\lambda_n)$ of distinct positive real
numbers, let $B(\lambda)$ be the least support diameter of a probability
measure $\mu$ on $\mathbb R$ for which
\[
 \int e^{i\lambda_jt}\,d\mu(t)=0\qquad(1\leq j\leq n).
\]
All Fourier frequencies here are angular frequencies; there is no $2\pi$ in
the exponential. We normalize the mean frequency to $1$.

\clearpage
\begin{theorem}[Explicit Hessian and transverse rigidity]
Suppose
\[
 0<\lambda_1<\cdots<\lambda_n<2,\qquad
 \lambda_j+\lambda_{n+1-j}=2.
\]
Put $j^*=n+1-j$, $r_j=e^{i\pi\lambda_j}$, and
\[
 W(z)=\frac{\prod_{j=1}^n(z-r_j)}{\prod_{j=1}^n(1-r_j)}
      =\sum_{k=0}^n w_kz^k.
\]
Assume $w_k>0$ for every $0\leq k\leq n$. Then $B(\lambda)=n\pi$,
and the unique optimum anchored at zero has masses $w_k$ at $k\pi$.
For $v\in\mathbb R^n$ with $\sum_jv_j=0$, define the real polynomial
\begin{equation}\label{new:eq:S}
 S_v(z)=-\frac\pi2\,zW(z)
              \sum_{j=1}^n\frac{v_j+v_{j^*}}{z-r_j}
        =\sum_{k=0}^n w_ks_k(v)z^k.
\end{equation}
The expression is a polynomial, $s_0(v)=s_n(v)=0$, and
\begin{equation}\label{new:eq:Q}
 Q_\lambda(v)=\frac1\pi\sum_{0\leq k<\ell\leq n}
  w_kw_\ell\bigl(s_k(v)-s_\ell(v)\bigr)^2
  \int_0^\infty\frac{z^{k+\ell-1}}{W(z)^2}\,dz
\end{equation}
satisfies
\[
 B(\lambda+\varepsilon v)
       =n\pi+\varepsilon^2Q_\lambda(v)+O_{\lambda,v}(\varepsilon^3).
\]
Thus the Hessian of $B$ restricted to $\sum\lambda_j=n$ is $2Q_\lambda$.
Its kernel is exactly
\[
 \{v: v_j+v_{j^*}=0\text{ for every }j\}.
\]
Consequently its rank is $\lceil n/2\rceil-1$, and it is strictly positive
on every nonzero direction transverse to the mirror locus.
\end{theorem}

\begin{proof}
We spell out the analytic and positivity inputs, because an arbitrary
solution of the moment equations need not minimize width.

\emph{The analytic minimizing branch.}
The reflection-paired theorem of the article supplies a nonnegative
exponential-polynomial certificate on $[0,n\pi]$ whose zero set is exactly
$\{0,\pi,\ldots,n\pi\}$. Its endpoint zeros are simple, its interior zeros
are double with positive second derivatives, and it lies in the real span of
$\cos(\lambda_jt),\sin(\lambda_jt)$. This last statement follows directly
from the keyhole-residue formula in the article, not from an assumption on
unproved RH inputs.

Let $\eta$ be near $\lambda$. Seek weights and nodes
$w_k(\eta)$, $0=u_0(\eta)<\cdots<u_n(\eta)$ solving total mass $1$ and
all $n$ complex annihilation equations. There are $2n+1$ unknown real
parameters and equations. At the base, a homogeneous variation in the
Jacobian kernel has exponential transform
\[
 \delta G(s)=A(e^{\pi s})+sC(e^{\pi s}),
\]
where $A,C$ are real polynomials of degree at most $n$, $A(1)=0$, and
$C(0)=0$. Its vanishing at $i\lambda_j$, together with the conjugate of
its vanishing at $i(2-\lambda_j)$, gives
$2iC(r_j)=0$. Since all $r_j$ are distinct, $W$ divides $C$; since
$C(0)=0$ and $W(0)>0$, $C=0$. Similarly $W$ divides $A$, and
$A(1)=0$, $W(1)=1$ imply $A=0$. Positivity of the weights implies that
all node variations vanish. The Jacobian is therefore invertible.

The real analytic implicit function theorem supplies an analytic branch of
positive weights and ordered nodes. To construct its certificate explicitly,
let $Y(t)$ be the vector of the $n$ cosine and $n$ sine moment functions,
and let $J$ be the full Jacobian in coordinates
$(w_0,\ldots,w_n,u_1,\ldots,u_n)$. Solve
$J^{\mathsf T}(a_0,b)=-e_{\mathrm{last}}$. The weight-coordinate equations
give $a_0+b\cdot Y(u_k)=0$ at every node. Taking the weighted sum shows
$a_0=0$, since the total mass is one and all moments vanish. Therefore
$q(t)=b\cdot Y(t)$ vanishes at every node, has derivative zero at each
interior node, and satisfies $w_nq'(u_n)=-1$. Invertibility of $J$
makes this normalized certificate unique and analytic. At the base it
is the positive scalar multiple of the reflection certificate with this
endpoint normalization. The strict endpoint slopes, interior
curvatures, and positivity on the remaining compact subintervals persist.
Thus the certificate remains nonnegative on $[0,u_n(\eta)]$ and zero only
at those nodes. Every feasible measure supported in that interval is
supported on those nodes. The mass and moment equations determine its
weights uniquely; both endpoint masses are positive. Hence no narrower
feasible measure exists: any such measure could be translated into this
interval, preserving its zero Fourier moments. Thus $B(\eta)=u_n(\eta)$ is analytic and globally
minimal in a neighbourhood of the base.

\emph{A defect identity valid at this general base.}
For the mean-preserving path $\eta=\lambda+\varepsilon v$, put
\[
 G_\varepsilon(s)=\sum_{k=0}^n
                 w_k(\varepsilon)e^{u_k(\varepsilon)s},\qquad
 p_k(x,\varepsilon)=
 \frac{w_k(\varepsilon)e^{xu_k(\varepsilon)}}{G_\varepsilon(x)}.
\]
For sufficiently small $\varepsilon$, $G_\varepsilon$ has exactly the $n$
zeros $i\eta_j$ in $0<\operatorname{Im}s<2$, counted with multiplicity.
Indeed the base $G_0(s)=W(e^{\pi s})$ has precisely those zeros there.
For $|\operatorname{Re}s|>R$ the first or last term dominates uniformly
for small $\varepsilon$, since the extreme weights remain bounded below
and all node gaps stay bounded below. On the finite rectangle
$|\operatorname{Re}s|\leq R$, $0\leq\operatorname{Im}s\leq2$, the base
has no boundary zeros, so the argument principle or Rouche's theorem
preserves the zero count. The prescribed simple zeros already account for
all $n$ zeros.

The article's horizontal-strip identity, at height $2$, is
\[
 \int_{\mathbb R}\log\frac{G_\varepsilon(x)}
             {|G_\varepsilon(x+2i)|}\,dx
  =2B(\eta)-2\pi\sum_{j=1}^n(2-\eta_j).
\]
Because $\sum\eta_j=n$, this becomes
\begin{equation}\label{new:eq:defect}
 B(\eta)-n\pi
 =\frac12\int_{\mathbb R}
   -\log\left|\sum_{k=0}^np_k(x,\varepsilon)
                       e^{2iu_k(\varepsilon)}\right|\,dx.
\end{equation}
This is precisely the general-base version of the defect formula used by
the article at the coalesced spectrum.

Let $s_k=\dot u_k(0)$. The elementary expansion of a characteristic
function at zero gives, for each fixed $x$,
\[
 -\log\left|\sum_kp_k(x,\varepsilon)e^{2iu_k(\varepsilon)}\right|
 =2\varepsilon^2\operatorname{Var}_{p(x,0)}(s_k)
      +O_x(\varepsilon^3).
\]
The variation of the weights does not affect the quadratic coefficient:
the phases agree at the base, and the weights always sum to one. The
exchange of differentiation and integration is justified as follows.
On a fixed compact $x$-interval the expression inside the logarithm is
uniformly nonzero for small $\varepsilon$. On either tail, factor out
the phase of the dominant endpoint; the sum of all other probabilities,
and its first three parameter derivatives, are bounded by
$C(1+|x|)^3e^{-g|x|}$ for fixed $g>0$. The logarithm and its first three
parameter derivatives have the same integrable bound. Taylor's theorem
under the integral therefore yields
\[
 B(\lambda+\varepsilon v)-n\pi
 =\varepsilon^2\int_{\mathbb R}
            \operatorname{Var}_{p(x,0)}(s_k)\,dx+O(\varepsilon^3).
\]
Now $p_k(x,0)=w_ke^{k\pi x}/W(e^{\pi x})$. Using
$\operatorname{Var}_p(s)=\sum_{k<\ell}p_kp_\ell(s_k-s_\ell)^2$ and
$z=e^{\pi x}$ gives exactly \eqref{new:eq:Q}. Every integral is finite
and strictly positive: near zero its power is at least zero, and near
infinity it is at most $z^{-2}$, since $k+\ell\leq2n-1$.

\emph{Explicit node velocities.}
Write
$A(z)=\sum_k\dot w_k(0)z^k$ and
$S(z)=\sum_kw_ks_kz^k$. Differentiating the $j$th annihilation equation
gives
\[
 A(r_j)+i\lambda_jS(r_j)
       +i\pi v_jr_jW'(r_j)=0.
\]
Conjugating the equation for $j^*$ and subtracting it gives
\[
 2iS(r_j)+i\pi(v_j+v_{j^*})r_jW'(r_j)=0.
\]
The polynomial in \eqref{new:eq:S} has exactly these values at every
$r_j$ and vanishes at $0$. These $n+1$ interpolation conditions determine
a polynomial of degree at most $n$ uniquely. Thus it is $S$. Its realness
also follows from conjugate pairing. Its coefficient of degree $n$ is
$-(\pi/2)w_n\sum_j(v_j+v_{j^*})=0$, and its constant coefficient is zero.
This proves $s_0=s_n=0$ and the velocity formula, without differentiating
the implicit optimizer numerically.

\emph{Kernel and rank.}
All edge weights in \eqref{new:eq:Q} are strictly positive. Thus $Q(v)=0$
if and only if every $s_k$ is equal. Since $s_0=0$, this means $S_v=0$.
The partial fraction expression in \eqref{new:eq:S}, whose poles are
distinct, vanishes identically if and only if $v_j+v_{j^*}=0$ for all $j$.
The mean-zero space has dimension $n-1$ and the reflection-preserving
subspace has dimension $\lfloor n/2\rfloor$, giving the stated rank.
\end{proof}

\begin{corollary}[Uniform quadratic distance from symmetry]
Near any base of the theorem, on the ordered mean-one hyperplane, put
\[
 m_j(\eta)=1+\frac{\eta_j-\eta_{j^*}}2,\qquad
 a_j(\eta)=\frac{\eta_j+\eta_{j^*}}2-1.
\]
Then $m$ is the orthogonal projection onto the local mirror locus and
$a=\eta-m$. There are constants $0<c<C$ and a neighbourhood on which
\[
 c\sum_ja_j(\eta)^2
 \ \leq\ B(\eta)-n\pi\ \leq\ C\sum_ja_j(\eta)^2.
\]
In particular the normalized width has a nondegenerate quadratic minimum
transverse to its entire local equality manifold.
\end{corollary}

\begin{proof}
The nearby mirror points $m$ remain distinct with every coefficient of
$W_m$ positive. The theorem applies uniformly on a compact neighbourhood
of these mirror points. Its Hessian restricted to the fixed normal
subspace is positive definite and depends continuously on $m$, hence its
least eigenvalue is bounded below by a positive constant. Expand
$B(m+a)$ at $m$ in the normal direction. The constant term is $n\pi$,
the linear term vanishes, the quadratic term is $Q_m(a)$, and the cubic
remainder is uniform. Shrinking the neighbourhood proves both bounds.
\end{proof}

\subsection{An exact noncoalesced three-frequency example}
Take
\[
 \lambda=(1/2,1,3/2),\quad W(z)=(1+z+z^2+z^3)/4,
 \quad v=(1,-2,1).
\]
Then
\[
 S_v(z)=\frac\pi2(z-z^2),\qquad
 (s_0,s_1,s_2,s_3)=(0,2\pi,-2\pi,0).
\]
Let $P(z)=1+z+z^2+z^3$. For $f=(0,1,-1,0)$ the variance integral is
\[
 \int_{\mathbb R}\operatorname{Var}_{p(x,0)}f\,dx
 =\frac1\pi\int_0^\infty
  \frac{(z+z^2)P(z)-(z-z^2)^2}{zP(z)^2}\,dz
 =\frac1\pi\int_0^\infty
  \left(\frac1{(z+1)^2}+\frac{z}{(z^2+1)^2}\right)\,dz
 =\frac3{2\pi}.
\]
Consequently
\[
 B(1/2+\varepsilon,1-2\varepsilon,3/2+\varepsilon)
 =3\pi+6\pi\varepsilon^2+O(\varepsilon^3).
\]
For any mean-zero $v\in\mathbb R^3$, the normal part is
$-(v_2/2)(1,-2,1)$, so the full restricted Hessian is described by
\[
 Q_\lambda(v)=\frac{3\pi}{2}v_2^2,
 \qquad D^2B(\lambda)[v,v]=3\pi v_2^2.
\]
The null direction $(1,0,-1)$ moves along the mirror locus.

More generally, at $\lambda=(1-d,1,1+d)$ with $0<d<2/3$, let
$c=\cos(\pi d)>-1/2$. Then
\[
 W(z)=\frac{(z+1)(z^2+2cz+1)}{4(1+c)},\quad
 S_{(1,-2,1)}(z)=\frac{\pi(1-c)}{2(1+c)}(z-z^2),
\]
and
\[
 s_1=-s_2=\frac{2\pi(1-c)}{1+2c},\qquad s_0=s_3=0.
\]
In fact the entire three-frequency transverse curvature has the closed
form
\begin{equation}\label{new:eq:threecurv}
 Q_{(1-d,1,1+d)}(1,-2,1)
 =6\pi\left(1-\frac{3\pi d\cos(\pi d)}
                    {(1+2\cos(\pi d))\sin(\pi d)}\right),
 \qquad 0<d<2/3.
\end{equation}
This positive expression has both a coalescence asymptotic and a
coefficient-boundary blow-up:
\[
 Q_{(1-d,1,1+d)}(1,-2,1)
   =\frac{3\pi^5}{10}d^4+\frac{3\pi^7}{28}d^6+O(d^8)
       \quad(d\downarrow0),
\]
\[
 Q_{(1-d,1,1+d)}(1,-2,1)
   \sim\frac{4\pi}{2/3-d}\quad(d\uparrow2/3).
\]
The leading $d^4$ term agrees with the article's sixth-order coalescence
law after differentiation in the symmetry-breaking direction; the displayed
$d^6$ term is a further term in the curvature expansion. The second asymptotic
quantifies the loss of smooth uniform control as the two interior masses
approach zero.

Here is an exact evaluation of the integral. Set $A=1+2c$ and
$P(z)=(z+1)(z^2+2cz+1)$. With $f=(0,1,-1,0)$ the variance density after
$z=e^{\pi x}$ is $R(z)\,dz/\pi$, where
\[
 R(z)=\frac{A(z+z^2)P(z)-A^2(z-z^2)^2}{zP(z)^2}.
\]
Partial fractions give the rational identity
\[
 R(z)=\frac{c(c-4)A}{(c-1)^2(z^2+2cz+1)}
      -\frac{z(c+1)A^2}{(c-1)(z^2+2cz+1)^2}
      +\frac{A^2}{(c-1)^2(z+1)^2}.
\]
Write $J(c)=\arccos(c)/\sqrt{1-c^2}$. Direct integration yields
\[
 \int_0^\infty\frac{dz}{z^2+2cz+1}=J(c),\qquad
 \int_0^\infty\frac{z\,dz}{(z^2+2cz+1)^2}
        =\frac{1-cJ(c)}{2(1-c^2)}.
\]
Therefore
\[
 \int_0^\infty R(z)\,dz
    =\frac{3A}{2(1-c)^2}\bigl(A-3cJ(c)\bigr).
\]
Multiplying by $s_1^2/\pi$ gives \eqref{new:eq:threecurv}. Taylor
expansion at $d=0$ gives the first asymptotic. At $d=2/3$,
$1+2\cos(\pi d)\sim\pi\sqrt3(2/3-d)$, while the remaining factors
have finite nonzero limits, giving the second. This calculation also
shows explicitly why the theorem excludes the zero-weight boundary.

\subsection{Novelty and limits}
The explicit velocity polynomial, all-base variance formula, exact kernel,
and quadratic-distance theorem extend the inspected article's real,
simple-spectrum mirror problem. Their proof develops its strip-defect
method. Established positive-input-shaping work, including
Crain--Singhose--Seering \cite{Crain1996}, provides related optimization
context; the accessible abstract discusses two-mode designs and their
mode-ratio dependence. Targeted primary-source searches did not locate
this arbitrary-$n$ Hessian formula. This is a new extension relative to
the inspected article; publication priority remains unverified.

For exponent targets $\beta+i\lambda_j$ with a common real $\beta$, the
positive tilt $d\nu=e^{\beta t}d\mu/\int e^{\beta u}d\mu$ preserves the
support and identifies the problem with the Fourier problem; all width
formulas above therefore apply unchanged. RH/Dirichlet GRH is needed only
if these targets are asserted to be zero contributions in a prime explicit
formula. This new theorem itself is unconditional and contains no
assumption about actual zero spacings, simplicity, or their independence.

\section{Verification status}
The proof uses the reflection certificate, analytic perturbation setup,
and strip identity of the source article, with the needed nonsingularity
and differentiation arguments given above. The new polynomial velocity
formula and the Hessian kernel do not require assumptions about zeta
zeros. The accompanying script checks the rational integral identities
symbolically and compares nine high-precision solutions of the moment
equations with the curvature formula. The numerical comparisons are
illustrations, not substitutes for the proof or interval certificates.
This companion has not been formalized or compiled in Lean.
{\small\raggedright
\begin{thebibliography}{9}
\bibitem{ProveItFilters} V. Reshetnikov, \emph{ProveIt: Beyond the Riemann
hypothesis for Dirichlet $L$-functions}, pinned revision
\texttt{cfab41369737c66ae520d179419509eb719367cd}, positive-filter section,
especially \texttt{fil:thm:reflect}, \texttt{fil:thm:perturb},
\texttt{fil:lem:strip}, \texttt{fil:thm:hessian}, and the second item of
\texttt{fil:sec:open}.\newline
\url{https://github.com/VladimirReshetnikov/ProveIt/blob/cfab41369737c66ae520d179419509eb719367cd/Analysis/GeneralizedRiemannHypothesis/Article/sections/fil.tex}.
\bibitem{Crain1996} E. A. Crain, W. E. Singhose, and W. P. Seering,
\emph{Derivation and properties of convolved and simultaneous two-mode
input shapers}, IFAC Proceedings Volumes 29 (1996), 7278--7283.
\url{https://doi.org/10.1016/S1474-6670(17)58856-4}.
\end{thebibliography}
}
\end{document}
